Zwischen Kopf und Füssen steht in den Blutgefässen eine zusammenhängende Blutsäule. Ihre Höhe ist etwa die Körpergrösse, wir schätzen $h \approx \hKO$. Blut ist etwas dichter als Wasser, $\rho_\text{B} \approx \rhoBO$. Der Blutdruck nimmt nach unten um den Schweredruck dieser Säule zu:
\begin{align}
 \Delta p &= \dpBF \\
   &= \rhoB\cdot\gE\cdot\hK \\
   &= \dpB \approx \result{\dpBbP{0}{2}}
\end{align}
Diese Zunahme ist vergleichbar mit dem Überdruck, den das Herz selbst erzeugt; deshalb misst man den Blutdruck auf Herzhöhe (am Oberarm).
